Training steps within a run are autocorrelated and may be nonstationary. We
retain fixed-window summaries and descriptive coefficients, not step-iid
bootstrap intervals, first-versus-last-window tests, significance stars, or
claims that a multiplicity correction calibrates invalid p-values. Resampling
two arms independently is an unpaired bootstrap. Independence between runs does
not imply independence between steps. A block bootstrap would require a justified
dependence model and is not automatically reliable for a short evolving trace.

Seed/run-level inference requires actual independent observations, an identified
estimand, and an eligible matched comparison. Reconstructed vectors from a mean
and standard error cannot support an empirical seed bootstrap. Trustworthy mean,
SD, and sample count could support an analytic summary-statistic Welch test,
but only after task and cohort identity are resolved. An all-zero trace has
zero pooled variance: Cohen's $d=0/0$ is \emph{not estimable}, not zero with a
zero-width interval. One seed cannot estimate across-seed uncertainty; this does
not preclude separately justified conditional prompt-level inference.

\paragraph{Multiplicity and eligibility.}
The historical five-test, 20-test BH, and 38-test Bonferroni families were not
consistent. Primary hypotheses were not preregistered. We withdraw their
significance decisions rather than choosing a smaller denominator. Before any
future inference, a single versioned comparison registry must specify comparison
ID, cohort IDs, independent unit, estimand, eligibility, family ID, raw p-value,
adjustment, and adjusted p-value. Retrospective enumeration must be labeled as
such; missing comparisons do not silently reduce a prespecified family.
For example, the historical $t(4)=7.44$ gives a two-sided $p\simeq0.001743$ and
$38p\simeq0.06623$, not a family-wide significant result. This arithmetic check
does not establish that the underlying comparison is scientifically eligible.

\paragraph{Power illustration, not evidence.}
\texttt{TTestIndPower} uses a pooled equal-variance model, not Welch's degrees
of freedom. For five observations per arm, the equal-variance 80\%-power MDE
at unadjusted $\alpha=0.05$ is approximately $d=2.024$; using $0.05/5$ or
$0.05/38$ gives approximately $2.744$ or $3.778$. These are planning
illustrations, not observed-effect ``achieved power,'' not an endorsement of
any historical correction family, and not evidence for a measured effect.

\paragraph{Historical seed summaries.}
The TRL Qwen2.5-0.5B reference is labeled GSM8K in its baseline description
and arithmetic elsewhere. Task identity remains unresolved. Its reported
five-seed mean $0.734$ and SD $0.0703$ are retained as historical summaries
only; neither task label is selected as authoritative, and no reconstructed
seed bootstrap or cross-task comparison is asserted. The later five-seed
Qwen3-8B comparison against a fixed base is separate from the reviewed record.
